\scanpage{13}

On en déduit des iso. can.

\[
\left\{\begin{aligned}
E^\vee&\simeq E\otimes L^\vee\\
\underbrace{\mathrm{Sym}^2(E^\vee)}_{\underline{\mathrm{Quad}}(E)}
&\simeq\underbrace{\mathrm{Sym}^2(E)}_{\underline{\mathrm{Quad}}(E^\vee)}\otimes
\underbrace{L^{\vee\otimes2}}_{L^{\otimes(-2)}}
\end{aligned}\right.
\]

\margin{i.e. $\underline{\mathrm{Quad}}(E)\simeq\mathrm{Quad}(E^\vee)\otimes L^{\otimes\unsure{-2}}$}

Il en résulte un isom. canonique de faisceaux

\[\underline{\mathrm{Sim}}(E)\xrightarrow{\substack{\mathrm{can}_E\\\sim}}\underline{\mathrm{Sim}}(E^\vee)\]

qui transforme le sous-Module loc. monogène $\Delta$ de $\underline{\mathrm{Sym}}^2(E^\vee)\simeq\mathrm{Quad}(E)$ dans

\[
\begin{aligned}
\Delta\otimes L^{\otimes(\unsure{2|-2})}&\subset
\deleted{\underline{\mathrm{Sym}}^2(E)}\ \underline{\mathrm{Quad}}(E^{\deleted{\unsure{\vee}}})\otimes L^{\otimes2}
\simeq \mathrm{Quad}(E^\vee)\otimes L^{\otimes(-2)}\inserted{\unsure{\otimes L^{\otimes2}|\otimes2}}\\
&\simeq\mathrm{Quad}(E^\vee).
\end{aligned}
\]

On trouve que $\mathrm{can}_E$ et $\mathrm{can}_{E^\vee}$ sont inverses l'un de l'autre. Les notions de \inserted{structures de} similitudes \uline{fidèles}, strictement fidèles, non-dégénérées (transversalement, universellement \ldots) \inserted{pour $E$ et pour $E^\vee$} se correspondent.

2.2. Considérons encore l'iso. can. (dépendant du choix d'un signe)

\[E^\vee\simeq E\otimes L^\vee,\]

en \inserted{$\otimes$-}multipliant \inserted{à droite} par $E$, on trouve

\[E^\vee\otimes E\simeq E\otimes E\otimes L^\vee\]

D'autre part on a \deleted{\unsure{un is. canonique}} \inserted{(pour tout $\underline{\mathcal O}_X$-Module)} une suite exacte canonique (qu'on peut mettre un $0$ à g si $E$ loc. Libre \ldots)

\[
\begin{aligned}
\bigwedge^2 E&\longrightarrow E\otimes E\longrightarrow\mathrm{Sym}^2(E)\longrightarrow0,\\
x\wedge y&\longrightarrow x\otimes y-y\otimes x
\end{aligned}
\]
% archive sheet number: 12
